\documentclass[10pt,a4paper]{amsart}
\usepackage{amsmath,amssymb,amsthm,graphicx}
% Retained preamble for the Stacks Project chapter snapshot in this directory.
% The chapter's own preamble.tex is a copy of the shared Stacks Project preamble; the block
% below is sliced MECHANICALLY, line for line, from it. Nothing was added, removed or
% redefined; the wrapper class, the xr-hyper/hyperref wiring and the externaldocument list
% are omitted because the assembler supplies the document class.
%
% lines taken from preamble.tex: 3, 5, 9, 11, 12, 156, 157, 158, ... (51 lines total)
%
% package coverage restored from the chapter's own preamble (2026-10-13): xy, xy-2cell

\usepackage[all]{xy}
\xyoption{2cell}
\UseAllTwocells

\newenvironment{reference}{\comment}{\endcomment}
\newenvironment{slogan}{\comment}{\endcomment}
\newenvironment{history}{\comment}{\endcomment}
\theoremstyle{plain}
\newtheorem{theorem}[subsection]{Theorem}
\newtheorem{proposition}[subsection]{Proposition}
\newtheorem{lemma}[subsection]{Lemma}
\theoremstyle{definition}
\newtheorem{definition}[subsection]{Definition}
\newtheorem{example}[subsection]{Example}
\newtheorem{exercise}[subsection]{Exercise}
\newtheorem{situation}[subsection]{Situation}
\theoremstyle{remark}
\newtheorem{remark}[subsection]{Remark}
\newtheorem{remarks}[subsection]{Remarks}
\numberwithin{equation}{subsection}
\def\lim{\mathop{\mathrm{lim}}\nolimits}
\def\colim{\mathop{\mathrm{colim}}\nolimits}
\def\Spec{\mathop{\mathrm{Spec}}}
\def\Hom{\mathop{\mathrm{Hom}}\nolimits}
\def\Ext{\mathop{\mathrm{Ext}}\nolimits}
\def\SheafHom{\mathop{\mathcal{H}\!\mathit{om}}\nolimits}
\def\SheafExt{\mathop{\mathcal{E}\!\mathit{xt}}\nolimits}
\def\Sch{\mathit{Sch}}
\def\Mor{\mathop{\mathrm{Mor}}\nolimits}
\def\Ob{\mathop{\mathrm{Ob}}\nolimits}
\def\Sh{\mathop{\mathit{Sh}}\nolimits}
\def\NL{\mathop{N\!L}\nolimits}
\def\CH{\mathop{\mathrm{CH}}\nolimits}
\def\proetale{{pro\text{-}\acute{e}tale}}
\def\etale{{\acute{e}tale}}
\def\QCoh{\mathit{QCoh}}
\def\Ker{\mathop{\mathrm{Ker}}}
\def\Im{\mathop{\mathrm{Im}}}
\def\Coker{\mathop{\mathrm{Coker}}}
\def\Coim{\mathop{\mathrm{Coim}}}
\DeclareMathSymbol{\boxtimes}{\mathbin}{AMSa}{"02}
\def\QCohstack{\mathcal{QC}\!\mathit{oh}}
\def\Cohstack{\mathcal{C}\!\mathit{oh}}
\def\Spacesstack{\mathcal{S}\!\mathit{paces}}
\def\Quotfunctor{\mathrm{Quot}}
\def\Hilbfunctor{\mathrm{Hilb}}
\def\Curvesstack{\mathcal{C}\!\mathit{urves}}
\def\Polarizedstack{\mathcal{P}\!\mathit{olarized}}
\def\Complexesstack{\mathcal{C}\!\mathit{omplexes}}
\def\Pic{\mathop{\mathrm{Pic}}\nolimits}
\def\Picardstack{\mathcal{P}\!\mathit{ic}}
\def\Picardfunctor{\mathrm{Pic}}
\def\Deformationcategory{\mathcal{D}\!\mathit{ef}}
\usepackage{cleveref}
\crefname{originalreference}{original reference}{original references}
\Crefname{originalreference}{Original reference}{Original references}
\makeatletter
\newlabel{morphisms-lemma-diagonal-unramified-morphism}{{02GE}{source}{}{sourcepointer-1}{}}
\newlabel{morphisms-lemma-diagonal-unramified-morphism@cref}{{[originalreference][0][]02GE}{[1][0][]source}{}{}{}}
\makeatother
\crefname{originalreference}{original reference}{original references}
\Crefname{originalreference}{Original reference}{Original references}
\makeatletter
\newlabel{more-morphisms-lemma-unramified-formally-unramified}{{02HE}{source}{}{sourcepointer-2}{}}
\newlabel{more-morphisms-lemma-unramified-formally-unramified@cref}{{[originalreference][0][]02HE}{[1][0][]source}{}{}{}}
\makeatother
\crefname{originalreference}{original reference}{original references}
\Crefname{originalreference}{Original reference}{Original references}
\makeatletter
\newlabel{more-morphisms-definition-universal-thickening}{{04F4}{source}{}{sourcepointer-3}{}}
\newlabel{more-morphisms-definition-universal-thickening@cref}{{[originalreference][0][]04F4}{[1][0][]source}{}{}{}}
\makeatother
\crefname{originalreference}{original reference}{original references}
\Crefname{originalreference}{Original reference}{Original references}
\makeatletter
\newlabel{more-morphisms-lemma-normalization-and-smooth}{{07TD}{source}{}{sourcepointer-4}{}}
\newlabel{more-morphisms-lemma-normalization-and-smooth@cref}{{[originalreference][0][]07TD}{[1][0][]source}{}{}{}}
\makeatother
\crefname{originalreference}{original reference}{original references}
\Crefname{originalreference}{Original reference}{Original references}
\makeatletter
\newlabel{more-algebra-lemma-regular-etale-extension}{{0AH0}{source}{}{sourcepointer-5}{}}
\newlabel{more-algebra-lemma-regular-etale-extension@cref}{{[originalreference][0][]0AH0}{[1][0][]source}{}{}{}}
\makeatother
\crefname{originalreference}{original reference}{original references}
\Crefname{originalreference}{Original reference}{Original references}
\makeatletter
\newlabel{descent-lemma-descending-property-unramified}{{02VM}{source}{}{sourcepointer-6}{}}
\newlabel{descent-lemma-descending-property-unramified@cref}{{[originalreference][0][]02VM}{[1][0][]source}{}{}{}}
\makeatother
\crefname{originalreference}{original reference}{original references}
\Crefname{originalreference}{Original reference}{Original references}
\makeatletter
\newlabel{more-morphisms-lemma-descending-property-lci}{{069N}{source}{}{sourcepointer-7}{}}
\newlabel{more-morphisms-lemma-descending-property-lci@cref}{{[originalreference][0][]069N}{[1][0][]source}{}{}{}}
\makeatother
\crefname{originalreference}{original reference}{original references}
\Crefname{originalreference}{Original reference}{Original references}
\makeatletter
\newlabel{descent-lemma-finite-locally-free-descends}{{05B2}{source}{}{sourcepointer-8}{}}
\newlabel{descent-lemma-finite-locally-free-descends@cref}{{[originalreference][0][]05B2}{[1][0][]source}{}{}{}}
\makeatother
\crefname{originallemma}{original lemma}{original lemmas}
\Crefname{originallemma}{Original lemma}{Original lemmas}
\makeatletter
\newlabel{lemma-strict-normal-crossings}{{0BIA}{source}{}{sourcepointer-9}{}}
\newlabel{lemma-strict-normal-crossings@cref}{{[originallemma][0][]0BIA}{[1][0][]source}{}{}{}}
\makeatother
\crefname{originalreference}{original reference}{original references}
\Crefname{originalreference}{Original reference}{Original references}
\makeatletter
\newlabel{more-morphisms-lemma-regular-immersion-lci}{{069L}{source}{}{sourcepointer-10}{}}
\newlabel{more-morphisms-lemma-regular-immersion-lci@cref}{{[originalreference][0][]069L}{[1][0][]source}{}{}{}}
\makeatother
\crefname{originalreference}{original reference}{original references}
\Crefname{originalreference}{Original reference}{Original references}
\makeatletter
\newlabel{divisors-lemma-Noetherian-scheme-regular-ideal}{{063I}{source}{}{sourcepointer-11}{}}
\newlabel{divisors-lemma-Noetherian-scheme-regular-ideal@cref}{{[originalreference][0][]063I}{[1][0][]source}{}{}{}}
\makeatother
\crefname{originalreference}{original reference}{original references}
\Crefname{originalreference}{Original reference}{Original references}
\makeatletter
\newlabel{morphisms-lemma-finite-integral}{{01WJ}{source}{}{sourcepointer-12}{}}
\newlabel{morphisms-lemma-finite-integral@cref}{{[originalreference][0][]01WJ}{[1][0][]source}{}{}{}}
\makeatother
\crefname{originallemma}{original lemma}{original lemmas}
\Crefname{originallemma}{Original lemma}{Original lemmas}
\makeatletter
\newlabel{lemma-finite-unramified-etale-local}{{04HJ}{source}{}{sourcepointer-13}{}}
\newlabel{lemma-finite-unramified-etale-local@cref}{{[originallemma][0][]04HJ}{[1][0][]source}{}{}{}}
\makeatother
\crefname{originalreference}{original reference}{original references}
\Crefname{originalreference}{Original reference}{Original references}
\makeatletter
\newlabel{more-morphisms-lemma-lci}{{069G}{source}{}{sourcepointer-14}{}}
\newlabel{more-morphisms-lemma-lci@cref}{{[originalreference][0][]069G}{[1][0][]source}{}{}{}}
\makeatother
\crefname{originalreference}{original reference}{original references}
\Crefname{originalreference}{Original reference}{Original references}
\makeatletter
\newlabel{divisors-lemma-regular-immersion-noetherian}{{063L}{source}{}{sourcepointer-15}{}}
\newlabel{divisors-lemma-regular-immersion-noetherian@cref}{{[originalreference][0][]063L}{[1][0][]source}{}{}{}}
\makeatother
\crefname{originalreference}{original reference}{original references}
\Crefname{originalreference}{Original reference}{Original references}
\makeatletter
\newlabel{algebra-lemma-regular-mod-x}{{00NU}{source}{}{sourcepointer-16}{}}
\newlabel{algebra-lemma-regular-mod-x@cref}{{[originalreference][0][]00NU}{[1][0][]source}{}{}{}}
\makeatother
\title{Normal crossings divisors characterized by the unramified normalization and regular fibre products}
\author{The Stacks Project Authors}
\date{stacks-project:etale.tex:local-snapshot-20261012; GNU FDL 1.2 or later}
\begin{document}
\maketitle
\noindent\textit{Editorial note.} All delivered mathematical text is the author's own, copied verbatim from the pinned Stacks Project chapter file etale.tex; no statement or proof step was rewritten, reordered or removed.\par\medskip
\noindent\textit{Editorial note.} Span recomputation by label pairing (LANE\_BRIEF 10.2): in this chapter the author's proof environment follows the statement environment as a standalone block; the spec covers the complete environment pair (statement span, proof span) as recorded in the pool candidate.\par\medskip
\ifdefined\hypertarget\hypertarget{sourcepointer-1}{}\fi
\noindent\textit{Source reference.} Original reference 02GE: sources/stacks-chow/tags.txt, source lines 3225--3225. Cross-chapter standard prerequisite cited by the author inside the counted span; it is used as a pointer and its own proof is not part of this item.\par\medskip
\ifdefined\hypertarget\hypertarget{sourcepointer-2}{}\fi
\noindent\textit{Source reference.} Original reference 02HE: sources/stacks-chow/tags.txt, source lines 3261--3261. Cross-chapter standard prerequisite cited by the author inside the counted span; it is used as a pointer and its own proof is not part of this item.\par\medskip
\ifdefined\hypertarget\hypertarget{sourcepointer-3}{}\fi
\noindent\textit{Source reference.} Original reference 04F4: sources/stacks-chow/tags.txt, source lines 5785--5785. Cross-chapter standard prerequisite cited by the author inside the counted span; it is used as a pointer and its own proof is not part of this item.\par\medskip
\ifdefined\hypertarget\hypertarget{sourcepointer-4}{}\fi
\noindent\textit{Source reference.} Original reference 07TD: sources/stacks-chow/tags.txt, source lines 9933--9933. Cross-chapter standard prerequisite cited by the author inside the counted span; it is used as a pointer and its own proof is not part of this item.\par\medskip
\ifdefined\hypertarget\hypertarget{sourcepointer-5}{}\fi
\noindent\textit{Source reference.} Original reference 0AH0: sources/stacks-chow/tags.txt, source lines 13213--13213. Cross-chapter standard prerequisite cited by the author inside the counted span; it is used as a pointer and its own proof is not part of this item.\par\medskip
\ifdefined\hypertarget\hypertarget{sourcepointer-6}{}\fi
\noindent\textit{Source reference.} Original reference 02VM: sources/stacks-chow/tags.txt, source lines 3773--3773. Cross-chapter standard prerequisite cited by the author inside the counted span; it is used as a pointer and its own proof is not part of this item.\par\medskip
\ifdefined\hypertarget\hypertarget{sourcepointer-7}{}\fi
\noindent\textit{Source reference.} Original reference 069N: sources/stacks-chow/tags.txt, source lines 8047--8047. Cross-chapter standard prerequisite cited by the author inside the counted span; it is used as a pointer and its own proof is not part of this item.\par\medskip
\ifdefined\hypertarget\hypertarget{sourcepointer-8}{}\fi
\noindent\textit{Source reference.} Original reference 05B2: sources/stacks-chow/tags.txt, source lines 6870--6870. Cross-chapter standard prerequisite cited by the author inside the counted span; it is used as a pointer and its own proof is not part of this item.\par\medskip
\ifdefined\hypertarget\hypertarget{sourcepointer-9}{}\fi
\noindent\textit{Source reference.} Original lemma 0BIA: sources/stacks-chow/tags.txt, source lines 14483--14483. Cross-chapter standard prerequisite cited by the author inside the counted span; it is used as a pointer and its own proof is not part of this item.\par\medskip
\ifdefined\hypertarget\hypertarget{sourcepointer-10}{}\fi
\noindent\textit{Source reference.} Original reference 069L: sources/stacks-chow/tags.txt, source lines 8045--8045. Cross-chapter standard prerequisite cited by the author inside the counted span; it is used as a pointer and its own proof is not part of this item.\par\medskip
\ifdefined\hypertarget\hypertarget{sourcepointer-11}{}\fi
\noindent\textit{Source reference.} Original reference 063I: sources/stacks-chow/tags.txt, source lines 7832--7832. Cross-chapter standard prerequisite cited by the author inside the counted span; it is used as a pointer and its own proof is not part of this item.\par\medskip
\ifdefined\hypertarget\hypertarget{sourcepointer-12}{}\fi
\noindent\textit{Source reference.} Original reference 01WJ: sources/stacks-chow/tags.txt, source lines 2484--2484. Cross-chapter standard prerequisite cited by the author inside the counted span; it is used as a pointer and its own proof is not part of this item.\par\medskip
\ifdefined\hypertarget\hypertarget{sourcepointer-13}{}\fi
\noindent\textit{Source reference.} Original lemma 04HJ: sources/stacks-chow/tags.txt, source lines 5870--5870. Cross-chapter standard prerequisite cited by the author inside the counted span; it is used as a pointer and its own proof is not part of this item.\par\medskip
\ifdefined\hypertarget\hypertarget{sourcepointer-14}{}\fi
\noindent\textit{Source reference.} Original reference 069G: sources/stacks-chow/tags.txt, source lines 8040--8040. Cross-chapter standard prerequisite cited by the author inside the counted span; it is used as a pointer and its own proof is not part of this item.\par\medskip
\ifdefined\hypertarget\hypertarget{sourcepointer-15}{}\fi
\noindent\textit{Source reference.} Original reference 063L: sources/stacks-chow/tags.txt, source lines 7835--7835. Cross-chapter standard prerequisite cited by the author inside the counted span; it is used as a pointer and its own proof is not part of this item.\par\medskip
\ifdefined\hypertarget\hypertarget{sourcepointer-16}{}\fi
\noindent\textit{Source reference.} Original reference 00NU: sources/stacks-chow/tags.txt, source lines 872--872. Cross-chapter standard prerequisite cited by the author inside the counted span; it is used as a pointer and its own proof is not part of this item.\par\medskip
\begin{lemma}
\label{lemma-characterize-normal-crossings-normalization}
Let $X$ be a locally Noetherian scheme. Let $D \subset X$ be a closed
subscheme. The following are equivalent
\begin{enumerate}
\item $D$ is a normal crossings divisor in $X$,
\item $D$ is reduced, the normalization $\nu : D^\nu \to D$ is unramified,
and for any $n \geq 1$ the scheme
$$
Z_n = D^\nu \times_D \ldots \times_D D^\nu
\setminus \{(p_1, \ldots, p_n) \mid p_i = p_j\text{ for some }i\not = j\}
$$
is regular, the morphism $Z_n \to X$ is a local complete intersection
morphism whose conormal sheaf is locally free of rank $n$.
\end{enumerate}
\end{lemma}

\begin{proof}
First we explain how to think about condition (2).
The diagonal of an unramified morphism is open
(Morphisms, Lemma \ref{morphisms-lemma-diagonal-unramified-morphism}).
On the other hand $D^\nu \to D$ is separated, hence the
diagonal $D^\nu \to D^\nu \times_D D^\nu$ is closed.
Thus $Z_n$ is an open and closed subscheme of
$D^\nu \times_D \ldots \times_D D^\nu$. On the other hand,
$Z_n \to X$ is unramified as it is the composition
$$
Z_n \to D^\nu \times_D \ldots \times_D D^\nu \to \ldots \to
D^\nu \times_D D^\nu \to D^\nu \to D \to X
$$
and each of the arrows is unramified.
Since an unramified morphism is formally unramified
(More on Morphisms, Lemma
\ref{more-morphisms-lemma-unramified-formally-unramified})
we have a conormal sheaf
$\mathcal{C}_n = \mathcal{C}_{Z_n/X}$ of $Z_n \to X$, see
More on Morphisms, Definition
\ref{more-morphisms-definition-universal-thickening}.

\medskip\noindent
Formation of normalization commutes with \'etale localization by
More on Morphisms, Lemma \ref{more-morphisms-lemma-normalization-and-smooth}.
Checking that local rings are regular, or that
a morphism is unramified, or that a morphism is a
local complete intersection or that a morphism is
unramified and has a conormal sheaf which is
locally free of a given rank, may be done \'etale locally (see
More on Algebra, Lemma \ref{more-algebra-lemma-regular-etale-extension},
Descent, Lemma \ref{descent-lemma-descending-property-unramified},
More on Morphisms, Lemma \ref{more-morphisms-lemma-descending-property-lci}
and
Descent, Lemma \ref{descent-lemma-finite-locally-free-descends}).

\medskip\noindent
By the remark of the preceding paragraph and the definition
of normal crossings divisor it suffices to prove that a
strict normal crossings divisor $D = \bigcup_{i \in I} D_i$
satisfies (2). In this case $D^\nu = \coprod D_i$
and $D^\nu \to D$ is unramified (being unramified
is local on the source and $D_i \to D$ is a closed
immersion which is unramified). Similarly, $Z_1 = D^\nu \to X$
is a local complete intersection morphism because we may
check this locally on the source and each morphism $D_i \to X$
is a regular immersion as it is the inclusion of a Cartier divisor
(see Lemma \ref{lemma-strict-normal-crossings} and
More on Morphisms, Lemma \ref{more-morphisms-lemma-regular-immersion-lci}).
Since an effective Cartier divisor has an invertible
conormal sheaf, we conclude that the requirement on the
conormal sheaf is satisfied.
Similarly, the scheme $Z_n$ for $n \geq 2$ is the disjoint union
of the schemes $D_J = \bigcap_{j \in J} D_j$ where $J \subset I$
runs over the subsets of order $n$. Since $D_J \to X$ is
a regular immersion of codimension $n$
(by the definition of strict normal crossings and the
fact that we may check this on stalks by
Divisors, Lemma \ref{divisors-lemma-Noetherian-scheme-regular-ideal})
it follows in the same manner that $Z_n \to X$ has the required
properties. Some details omitted.

\medskip\noindent
Assume (2). Let $p \in D$. Since $D^\nu \to D$ is unramified, it is
finite (by Morphisms, Lemma \ref{morphisms-lemma-finite-integral}).
Hence $D^\nu \to X$ is finite unramified.
By Lemma \ref{lemma-finite-unramified-etale-local}
and \'etale localization (permissible by the discussion
in the second paragraph and the definition of normal
crossings divisors) we reduce to the case where
$D^\nu = \coprod_{i \in I} D_i$
with $I$ finite and $D_i \to U$ a closed immersion.
After shrinking $X$ if necessary, we may assume
$p \in D_i$ for all $i \in I$. The condition that $Z_1 =  D^\nu \to X$ is an
unramified local complete intersection morphism
with conormal sheaf locally free of rank $1$
implies that $D_i \subset X$ is an effective Cartier divisor, see
More on Morphisms, Lemma \ref{more-morphisms-lemma-lci} and
Divisors, Lemma \ref{divisors-lemma-regular-immersion-noetherian}.
To finish the proof we may assume $X = \Spec(A)$ is affine
and $D_i = V(f_i)$ with $f_i \in A$ a nonzerodivisor.
If $I = \{1, \ldots, r\}$, then $p \in Z_r = V(f_1, \ldots, f_r)$.
The same reference as above implies that
$(f_1, \ldots, f_r)$ is a Koszul regular ideal in $A$.
Since the conormal sheaf has rank $r$, we see that
$f_1, \ldots, f_r$ is a minimal set of generators of
the ideal defining $Z_r$ in $\mathcal{O}_{X, p}$.
This implies that $f_1, \ldots, f_r$ is a regular sequence
in $\mathcal{O}_{X, p}$ such that $\mathcal{O}_{X, p}/(f_1, \ldots, f_r)$
is regular. Thus we conclude by
Algebra, Lemma \ref{algebra-lemma-regular-mod-x}
that $f_1, \ldots, f_r$ can be extended to a regular system of parameters
in $\mathcal{O}_{X, p}$ and this finishes the proof.
\end{proof}

\end{document}
